\section{Exact orbit counting from cycle inventories}
\label{sec:counting}
Let $a_G(\lambda)$ count the elements of $G$ whose cycle-length multiset is $\lambda$. Here \emph{fixed points are cycles of length one}. Define the cycle-index polynomial
\[
  Z_G(z_1,\ldots,z_n)=\frac1{|G|}
  \sum_{\lambda} a_G(\lambda)\prod_{\ell\in\lambda}z_\ell.
\]
The certificate stores the integer inventory $a_G$ and denominator $|G|$, not approximate rational coefficients. Its cycle partitions are sorted, multiplicities are positive integers, duplicate partitions are forbidden, and the multiplicities must sum to the independently checked group order.

Those syntactic checks are necessary but insufficient. The producer constructs the inventory by a Cayley-graph breadth-first enumeration from the original generators. The checker uses a genuinely different representation: it enumerates every unique transversal normal form from Theorem~\ref{thm:chain}, recomputes its cycles, and compares the entire inventory. It never calls the producer's group enumeration.

\subsection{A formula valid for every number of colors}
Burnside's lemma gives
\begin{equation}
  N_G(q)=\#([q]^{[n]}/G)
  =\frac1{|G|}\sum_{\lambda}a_G(\lambda)q^{|\lambda|}.
  \label{eq:burnside}
\end{equation}
An assignment fixed by a permutation is constant on each of its cycles, so a permutation with $k$ cycles has exactly $q^k$ fixed assignments. The averaging theorem is already available in Mathlib as
\begin{lstlisting}
MulAction.sum_card_fixedBy_eq_card_orbits_mul_card_group
\end{lstlisting}
in \code{Mathlib.GroupTheory.GroupAction.Quotient} \cite{mathlib-burnside}. This is a concrete library reuse point, not a request to refound finite group actions.

A future Lean implementation should prove the integer form
\[
 |G|\,N_G(q)=\sum_{\lambda}a_G(\lambda)q^{|\lambda|},
\]
then discharge exact division using $|G|>0$. The Python checker verifies divisibility for each evaluated $q$. The mathematical derivation establishes why one verified inventory supports a theorem quantified over $q$; executing finitely many values of $q$ alone would not establish that universal theorem.

\begin{example}[Ten positions on a circle]
Let $C_{10}$ act by rotation and $D_{10}$ by rotations and reflections of ten positions. The actions are explicit permutations; $|C_{10}|=10$ and $|D_{10}|=20$. Their orbit-count formulas are
\begin{align*}
 N_{C_{10}}(q)&=\frac{q^{10}+q^5+4q^2+4q}{10},\\
 N_{D_{10}}(q)&=\frac{q^{10}+5q^6+6q^5+4q^2+4q}{20}.
\end{align*}
For two colors these give 108 necklaces and 78 bracelets. For three colors the bracelet count is 3,210. The $D_{10}$ numerator is checked by the regression suite, rather than inserted as an unexplained formula.
\end{example}

\subsection{Fixed-weight counting}
For binary colorings with exactly $k$ ones, a fixed assignment selects whole cycles. Therefore
\begin{equation}
 N_{G,k}=\frac1{|G|}\sum_{\lambda}a_G(\lambda)
 [t^k]\prod_{\ell\in\lambda}(1+t^\ell).
 \label{eq:weight}
\end{equation}
The checker computes each truncated product by an integer dynamic program, updating coefficients in descending order so that a cycle is selected at most once. This detail matters: an ascending in-place update would accidentally permit multiple selections of the same cycle.

For $D_{10}$ and $k=5$, the result is 16. The main experiment compares all tested weights with an independent coloring-orbit oracle; weights above $n$ return zero. The empty action at $n=0$ has one empty coloring even when $q=0$, so the implementation treats $0^0=1$ in precisely that combinatorial context.

\subsection{Counting huge recognized groups without enumerating them}
For $n>0$,
\begin{equation}
 N_{\Sym_n}(q)=\binom{q+n-1}{n},
 \label{eq:symmetric-count}
\end{equation}
with value zero when $q=0$. Each orbit is a multiset of $n$ colors. The $n=0$ value is one.

For $n\geq2$,
\begin{equation}
 N_{\Alt_n}(q)=\binom{q+n-1}{n}+\binom qn.
 \label{eq:alternating-count}
\end{equation}
Here $\binom qn=0$ when $q<n$. To prove the additional term, examine a symmetric-group orbit. If a color repeats, its stabilizer contains an odd transposition, and the orbit remains one alternating-group orbit. If every color is distinct, its stabilizer is trivial and its symmetric-group orbit splits into two alternating-group orbits. There are $\binom qn$ such all-distinct color sets. This also explains the parity distinction in Theorem~\ref{thm:alternating-image}.

Thus a recognized $\Sym_{40}$ certificate immediately supports 41 binary-color orbits and 861 three-color orbits. No cycle inventory of $40!$ elements is constructed. The general inventory worker remains bounded by explicit group enumeration; the family formulas are a distinct certified route around that limitation.

\subsection{The singleton-cycle adapter hazard}
Mathlib's \code{Equiv.Perm.cycleType} omits fixed points: its entries are at least two and its sum is the support cardinality. This is visible in the documented theorems \code{two_le_of_mem_cycleType} and \code{sum_cycleType} \cite{mathlib-cycles}. Our cycle inventory includes fixed points. Consequently a direct translation using only the cardinality of \code{cycleType} would undercount cycles, including giving zero cycles for the identity.

The adapter must reconstruct the full multiset by adjoining $n-|\supp(g)|$ copies of one. Equivalently, the exponent in \eqref{eq:burnside} is
\[
  \#\operatorname{cycleType}(g)+n-|\supp(g)|.
\]
This is an example of a small representation mismatch that can invalidate every otherwise correct counting certificate.

\subsection{Constraints require a different fixed-point calculation}
Equation~\eqref{eq:burnside} counts all colorings, not satisfying assignments of an arbitrary CNF and not proper graph colorings. For a $G$-invariant constraint $P$, the correct numerator is
\[
  \sum_{g\in G}\#\{c:P(c)\ \text{and}\ g\cdot c=c\}.
\]
A proposed extension substitutes one variable per permutation cycle into $P$ and invokes a separate exact finite counting worker. The substitution's equivalence and each fixed-point count need certificates. This constrained-counting compiler is designed only; the delivered counting implementation covers unrestricted colors and fixed binary weight.
